\begin{proof}
We prove \cref{thm:p171-full-assembly}.  We first check that the source
hypotheses place the family in the exact
finite-presentation scope of the preceding algebraic results.
Pille--Schneider assume that \(A\) is integral and that \(X\) is projective
of finite type over \(A\), while their question adds flatness
\cite[Section~2 and Question~4.9]{PS23Hybrid}.  Flatness and local finite
type over an integral base imply local finite presentation by
\cite[Tag~05IN]{stacks-project}.  Finite type makes the
morphism quasi-compact, while projectivity makes it separated and hence
quasi-separated.  The definition in \cite[Tag~01TP]{stacks-project}
therefore makes
\(X\to\operatorname{Spec}A\) a morphism of finite presentation.  No
Noetherian or reducedness hypothesis has been added.

We separate the finite-level assertion from the density passage.  Fix a
projective embedding
\[
 \iota:X\hookrightarrow\mathbf P_A^r,
 \qquad H=\iota^*\mathcal O(1),
 \qquad \eta=\max_{0\leq k\leq r}\log|t_k|,
\]
where the \(t_k\) are the pulled-back coordinate sections.  Let
\(\phi_i^\circ\) be a tropical Fubini--Study metric on \(L_i\), and let
\(\psi^\circ\) be a DFS function.  By
\cite[Definition~2.15]{PS23Hybrid}, after fixing the isomorphism in the
DFS presentation,
write
\[
 \psi^\circ=\alpha-\beta
\]
with \(\alpha\) and \(\beta\) tropical Fubini--Study metrics on the same
semiample line bundle \(H_0\).
\Cref{prop:p171-closed-immersion-finite-level-reduction} gives
\begin{align}
 \zeta^+&=\alpha+\eta,&
 \zeta^-&=\beta+\eta,&
 \psi^\circ&=\zeta^+-\zeta^-,
 \label{eq:p171-full-dfs-expansion}\\
 \theta_i^+&=\phi_i^\circ+\eta,&
 \theta_i^-&=\eta,&
 \phi_i^\circ&=\theta_i^+-\theta_i^-.
 \label{eq:p171-full-fs-expansion}
\end{align}
Write \(G_i^+=L_i+H\) and \(G_i^-=H\) for the underlying line
bundles of \(\theta_i^+\) and \(\theta_i^-\), respectively.
Every displayed metric is defined by a finite system whose associated
morphism is a closed immersion: the construction is the Segre product of
the original basepoint-free morphism with a Veronese power of \(\iota\).
Consequently each underlying line bundle in the displayed positive terms
is ample.

For \(\varepsilon=(\varepsilon_1,\ldots,\varepsilon_n)\in
\{+,-\}^n\), set
\[
 s(\varepsilon)=(-1)^{\#\{i:\varepsilon_i=-\}},
 \qquad
 \boldsymbol\theta^\varepsilon
 =(\theta_1^{\varepsilon_1},\ldots,
   \theta_n^{\varepsilon_n}).
\]
Symmetry and multilinearity of the fieldwise mixed Monge--Amp\`ere
pairing give
\begin{equation}
 \operatorname{MA}(\phi_1^\circ,\ldots,\phi_n^\circ)
 =\sum_{\varepsilon\in\{+,-\}^n}s(\varepsilon)
   \operatorname{MA}(\boldsymbol\theta^\varepsilon).
 \label{eq:p171-full-ma-expansion}
\end{equation}
For a fixed \(\varepsilon\), the two metrics \(\zeta^+\) and
\(\zeta^-\) live on the same line bundle.  Let
\[
 \mathscr D_\varepsilon
 =I_{X/A}(H_0+H,G_1^{\varepsilon_1},\ldots,
   G_n^{\varepsilon_n})
\]
be the corresponding algebraic intersection line and choose a local frame
\(e_\varepsilon\).  The Deligne change-of-metric formula
\cite[Theorem~8.16, formula~(8.9)]{BE21} identifies their integral difference with
the difference of the two additive Deligne weights.  Thus, on any local
frame of the relevant intersection line,
\begin{align}
 &\int_{X_x^{\mathrm{an}}}\psi_x^\circ\,
     \operatorname{MA}_x
       ((\phi_1^\circ)_x,\ldots,(\phi_n^\circ)_x)\notag\\
 &\quad=
 \sum_{\varepsilon\in\{+,-\}^n}s(\varepsilon)
 \left(
 \lambda_x(e_\varepsilon(x);
           \zeta_x^+,\boldsymbol\theta_x^\varepsilon)
 -\lambda_x(e_\varepsilon(x);
           \zeta_x^-,\boldsymbol\theta_x^\varepsilon)
 \right).
 \label{eq:p171-full-signed-weights}
\end{align}
The right side consists literally of \(2\cdot2^n=2^{n+1}\) signed
Deligne weights.  If \(n=0\), the set \(\{+,-\}^0\) has one element,
so \eqref{eq:p171-full-signed-weights} has exactly two weights.  For each
fixed \(\varepsilon\) the two frame-change terms cancel, and hence their
difference is a globally defined real function.

We now justify continuity of every weight occurring in
\eqref{eq:p171-full-signed-weights}.  Consider one of its
closed-immersion tuples and write it in the notation of
\cref{thm:p171-module-c-synchronized-comparison}:
\[
 \vartheta_i=\frac1{m_i}\max_j
 \{\log|s_{ij}|+a_{ij}\},
 \qquad
 s_{ij}\in H^0(X,M_i^{m_i}),
 \qquad 0\leq i\leq n.
\]
Each \(M_i\) is ample and each displayed section system defines a closed
immersion.  Put
\begin{equation}
 d_i=(M_0\cdots\widehat M_i\cdots M_n)[X_x]_n,
 \qquad
 Q_p=\prod_{j=0}^n(m_jp),
 \qquad
 \delta_{i,p}=\frac{Q_p}{m_ip}d_i.
 \label{eq:p171-full-degree-ledger}
\end{equation}
Flatness makes the \(d_i\), hence the block degrees, independent of \(x\).
For \(n=0\), the same ledger reads
\(Q_p=m_0p\) and \(\delta_{0,p}=d_0\), where \(d_0\) is the degree of
\([X_x]_0\).

Put
\[
 \mathscr L=I_{X/A}(M_0,\ldots,M_n).
\]
By \cref{cor:p171-revised2-veronese-resultants}, for every sufficiently
large \(p\) there is a raw tensor
\[
 \mathcal R_p^{\mathrm{raw}}\in
 I_{X/A}(M_0^{m_0p},\ldots,M_n^{m_np})\otimes_A
 \bigotimes_{i=0}^n
 \operatorname{Sym}^{\delta_{i,p}}(E_{i,p}^\vee),
 \qquad
 E_{i,p}=\operatorname{Sym}^p(A^{N_i+1}).
\]
Ordered positive multiadditivity gives the isomorphism
\[
 \iota_p:
 I_{X/A}(M_0^{m_0p},\ldots,M_n^{m_np})
 \xrightarrow{\sim}\mathscr L^{\otimes Q_p}.
\]
Set
\[
 \mathcal R_p=(\iota_p\otimes1)
  (\mathcal R_p^{\mathrm{raw}})
 \in\mathscr L^{\otimes Q_p}\otimes_A
 \bigotimes_{i=0}^n
 \operatorname{Sym}^{\delta_{i,p}}(E_{i,p}^\vee).
\]
The construction uses one fixed Noetherian stage.  On every residue-field
fibre,
\(\mathcal R_p^{\mathrm{raw}}\) is the positive Chen--Moriwaki resultant
\cite[Section~1.6]{ChenMoriwaki21AdelicIntersection} of \([X_x]_n\),
including every top generic length, and \(\mathcal R_p\)
is its transport under \(\iota_p\).  Affine monomial
substitution is performed only after the complete-fibre theorem; kernels
in the substitution map therefore require neither an inherited norm open
nor a projective coefficient map.

Let \(x_0\in\mathcal M(A)\).  Because only finitely many intersection
lines occur in \eqref{eq:p171-full-signed-weights}, there is one common
analytic neighbourhood \(U\) of \(x_0\), independent of \(p\), on which
all of them have frames.  For each fixed \(p\) and each point of \(U\),
refine locally to a Banach analytic domain whose algebra contains the
finitely many tensor coefficients.  These domains cover \(U\), so a
coefficient weight proved continuous on each of them is continuous on
\(U\).  For a frame \(e\) of \(\mathscr L^{\mathrm{an}}\), write
\[
 \mathcal R_p=e^{\otimes Q_p}F_{p,e},
 \qquad
 F_{p,e}=\sum_\gamma c_{p,\gamma}\mathbf U^\gamma.
\]
The fibre resultant is nonzero at every point, so the coefficient tuple is
nowhere zero on \(U\).
\Cref{thm:hybrid-weighted-philippon-gauss} therefore makes the following
coefficient functional continuous.  With
\[
 \rho(x)=\frac{\log\max\{1,|2|_x\}}{\log2},
 \qquad
 w_{p,\gamma}=\sum_{i,\beta}
       \gamma_{i,\beta}\sum_j\beta_ja_{ij},
\]
it is the weighted Gauss maximum
\[
 \Phi_{p,e}(x)=
 \max_\gamma\{\log|c_{p,\gamma}(x)|_x+w_{p,\gamma}\},
 \qquad \log0:=-\infty,
\]
when \(\rho(x)=0\).  When \(\rho(x)>0\), it is
\begin{align*}
 \Phi_{p,e}(x)=\rho(x)\Bigg(&
 \int_{\prod_iS^{2K_{i,p}-1}}
 \log\left|
  \sum_\gamma\chi_x(c_{p,\gamma})
  \exp(w_{p,\gamma}/\rho(x))\mathbf U^\gamma
 \right|_\infty\,\mathrm d\eta\\
 &+\frac12\sum_i\delta_{i,p}
    \sum_{\ell=1}^{K_{i,p}-1}\frac1\ell\Bigg),
 \qquad K_{i,p}=\binom{N_i+p}{p}.
\end{align*}
The factor \(\rho(x)\) multiplies both the sphere integral and the
harmonic correction.  The representing character is chosen only on the
local Banach algebra and is unique up to conjugation.  The weighted theorem
proves choice independence and continuity through \(\rho=0\).  On an
overlap (e'=ue), its exact frame law is
\begin{equation}
 F_{p,e'}=u^{-Q_p}F_{p,e},
 \qquad
 \Phi_{p,e'}=\Phi_{p,e}-Q_p\log|u|.
 \label{eq:p171-full-frame-law}
\end{equation}

\Cref{lem:p171-module-c-field-comparison,thm:p171-module-c-synchronized-comparison}
prove on the
trivially valued, nontrivially non-Archimedean, powered real, and powered
complex fibres that
\begin{equation}
 \lambda_{p,e}(x)=Q_p^{-1}\Phi_{p,e}(x).
 \label{eq:p171-full-module-c-identity}
\end{equation}
Here the quotient metrics use the diagonal monomial basis and all powers in
\eqref{eq:p171-full-degree-ledger}.  The real measure is pushed forward
from complexification with no factor two.  In dimension zero the identity
is the positive finite norm, including residue degrees and Artin lengths.

The same theorem proves, without using continuity of the target weight,
that its normalized quotient metrics satisfy
\begin{equation}
 0\leq\vartheta_{i,p,x}-\vartheta_{i,x}
 \leq
 \frac{\rho(x)\log K_{i,p}}{2m_ip}
 \leq\frac{\log K_{i,p}}{2m_ip},
 \label{eq:p171-full-softmax}
\end{equation}
with equality in the first comparison at every non-Archimedean point.
The positively oriented change-of-metric formula
\cite[formula~(8.9)]{BE21} and the
degree identity
\(\delta_{i,p}(m_ip)=Q_pd_i\) then give
\begin{equation}
 0\leq\lambda_{p,e}(x)-\lambda_e(x)
 \leq E_p:=\sum_{i=0}^n
 d_i\frac{\log K_{i,p}}{2m_ip}
 \xrightarrow[p\to\infty]{}0.
 \label{eq:p171-full-global-error}
\end{equation}
The number \(E_p\) is independent of \(x\) and of the chosen local frame.
Thus the approximation is global even though its coefficients and frames
are read on local Banach neighbourhoods.  Since
\(\log K_{i,p}=O(\log p)\), it converges to zero.  Equations
\eqref{eq:p171-full-module-c-identity} and
\eqref{eq:p171-full-global-error} show that every target Deligne weight is
a local uniform limit of continuous coefficient weights.  It is therefore
continuous, and division of \eqref{eq:p171-full-frame-law} by \(Q_p\)
gives the correct gluing law.  This order avoids circularity: algebraic
tensor first, independently continuous coefficient functional second,
fieldwise metric identity third, and only then the uniform limit.

\Cref{cor:p171-module-c-finite-level-continuity} applies to
\(\psi^\circ,\phi_1^\circ,\ldots,\phi_n^\circ\).  Equivalently, combine
the \(2^{n+1}\) continuous weights in
\eqref{eq:p171-full-signed-weights}.  Hence
\begin{equation}
 x\longmapsto
 \int_{X_x^{\mathrm{an}}}\psi_x^\circ\,
 \operatorname{MA}_x
 ((\phi_1^\circ)_x,\ldots,(\phi_n^\circ)_x)
 \label{eq:p171-full-finite-level}
\end{equation}
is continuous for every DFS test and every tuple of tropical
Fubini--Study metrics.  This proves the finite-level continuity needed
below.

We now pass to arbitrary data by a cofinal product net.
Let
\[
 M=(L_1\cdots L_n)[X_x]_n\geq0
\]
be the common fibre mass; for \(n=0\), this is the degree of
\([X_x]_0\).  By \cite[Theorem~2.16]{PS23Hybrid}, there is a net of DFS functions
\(\psi_q\to\psi\) uniformly.  Write
\(\psi_q=\alpha_q-\beta_q\) on one semiample line bundle \(H_q\), using
the isomorphism in its DFS presentation.  By
\cite[Proposition~2.21]{PS23Hybrid}, independently for each \(i\), there is a net of tropical
Fubini--Study metrics
\(\phi_{i,a_i}\to\phi_i\) uniformly.

\Cref{cor:p171-production-family-CLN} applies to the slot-zero pair
\((\alpha_q,\beta_q)\); the arbitrary continuous function \(\psi\) is
never inserted into that CLN slot.  Put
\begin{equation}
 C_q=
 \begin{cases}
  0,&n=0,\\[1mm]
  2\displaystyle\max_{1\leq i\leq n}
   (H_qL_1\cdots\widehat L_i\cdots L_n)[X_x]_n,&n\geq1.
 \end{cases}
 \label{eq:p171-full-cln-constant}
\end{equation}
Flatness makes \(M\) and every \(C_q\) independent of \(x\).  At powered
Archimedean points, \eqref{eq:PN1-r3} leaves the unpowered mixed measure
unchanged, while \eqref{eq:PN2-r3} contributes one factor \(\rho\), and the
reciprocal
\(1/\rho\) from unpowering the varying metric difference cancels it.
At \(\rho=0\), \cite[Lemma~4.10]{PS23Hybrid} supplies the endpoint without
division by \(\rho\).  Hence the same \(C_q\) works on every fibre.  No
bound uniform in \(q\) is asserted or needed.

Let \(Q\) and \(A_i\) be the directed sets indexing these nets and put
\begin{equation}
 P=(0,1]^{\mathrm{op}}\times Q\times\prod_{i=1}^nA_i
 \label{eq:p171-full-product-set}
\end{equation}
with coordinatewise order.  Let \(D\subset P\) consist of
\(d=(\epsilon,q,\mathbf a)\) satisfying
\begin{align}
 (1+M)\|\psi_q-\psi\|_\infty&\leq\epsilon/2,
 \label{eq:p171-full-D-test}\\
 (1+C_q)\sum_{i=1}^n
 \|\phi_{i,a_i}-\phi_i\|_\infty&\leq\epsilon/2.
 \label{eq:p171-full-D-metric}
\end{align}
The set \(D\) is directed.  Given two of its elements, first choose a
smaller \(\epsilon\), next a common later \(q\) satisfying
\eqref{eq:p171-full-D-test}, and only after that \(q\) is fixed choose
common later \(a_i\)'s satisfying \eqref{eq:p171-full-D-metric}.

The cofinality is equally literal.  Given
\(p_0=(\epsilon_0,q_0,\mathbf a_0)\in P\), first choose
\(0<\epsilon\leq\epsilon_0\).  Uniform convergence of the DFS net gives
\(q\geq q_0\) satisfying \eqref{eq:p171-full-D-test}; after this
particular \(q\) and \(C_q\) are fixed, uniform convergence of the finitely
many metric nets gives \(a_i\geq a_{0,i}\) satisfying
\eqref{eq:p171-full-D-metric}.  Hence
\(d=(\epsilon,q,\mathbf a)\in D\) and \(d\geq p_0\).

For \(d\in D\), put
\begin{align*}
 F_d(x)&=\int_{X_x^{\mathrm{an}}}\psi_{q,x}\,
   \operatorname{MA}_x
   ((\phi_{1,a_1})_x,\ldots,(\phi_{n,a_n})_x),\\
 F(x)&=\mathcal I_{\psi,\boldsymbol\phi}(x).
\end{align*}
Each \(F_d\) is continuous by \eqref{eq:p171-full-finite-level}.  Positivity
and common mass control the test-function error, and
\cref{cor:p171-production-family-CLN}, applied only to
\(\alpha_q-\beta_q\), controls the metric error.  Uniformly for every
\(x\in\mathcal M(A)\),
\begin{align}
 |F_d(x)-F(x)|
 &\leq M\|\psi_q-\psi\|_\infty
 +C_q\sum_{i=1}^n
   \|\phi_{i,a_i}-\phi_i\|_\infty\notag\\
 &\leq\epsilon.
 \label{eq:p171-full-product-error}
\end{align}
Thus \(F_d\to F\) uniformly along the cofinal directed set \(D\).  A
uniform limit of a net of continuous real-valued functions is continuous
on an arbitrary topological space.  This proves
\cref{thm:p171-full-assembly}.
\end{proof}
